\documentclass[11pt]{article}
\usepackage[T1]{fontenc}
\usepackage[margin=1in]{geometry}
\usepackage{amsmath}
\usepackage[demo]{graphicx}
\usepackage{hyperref}
\hypersetup{colorlinks=true,linkcolor=blue,urlcolor=blue,citecolor=blue}

\title{FlashTeX Rendering Pipeline: Architecture Reference}
\author{FlashTeX Documentation Team}
\date{\today}

\begin{document}
\maketitle

\begin{abstract}
This document is an internal reference describing the stages of the
FlashTeX rendering pipeline, from source parsing through final PDF
emission. It is intended as a map of the codebase for new contributors and
exercises the cross-referencing machinery --- table of contents, list of
figures, list of tables, labels, and hyperlinks --- that the pipeline itself
must reproduce faithfully when compared against \texttt{pdflatex}.
\end{abstract}

\tableofcontents
\listoffigures
\listoftables

\section{Overview}
\label{sec:overview}

The pipeline is organized into four major stages, summarized in
Table~\ref{tab:stages} on page~\pageref{tab:stages}. Each stage is described
in its own section below, and Figure~\ref{fig:pipeline} on
page~\pageref{fig:pipeline} shows the data flow between them.

\begin{table}[ht]
  \centering
  \begin{tabular}{|l|l|l|}
    \hline
    Stage & Input & Output \\
    \hline
    Parsing & \texttt{.tex} source & Token stream \\
    Layout & Token stream & Box tree \\
    Shaping & Box tree & Glyph runs \\
    Emission & Glyph runs & PDF bytes \\
    \hline
  \end{tabular}
  \caption{The four pipeline stages}
  \label{tab:stages}
\end{table}

\begin{figure}[ht]
  \centering
  \includegraphics[draft,width=0.6\textwidth]{pipeline-diagram}
  \caption{High-level data flow through the rendering pipeline}
  \label{fig:pipeline}
\end{figure}

See also Section~\ref{sec:appendix} for a glossary of terms used throughout
this reference.\footnote{An expanded glossary, kept in sync with the
codebase, is published at \url{https://example.org/flashtex/glossary}.}

\section{Parsing}
\label{sec:parsing}

\subsection{Tokenization}
\label{sec:tokenization}

The tokenizer converts raw source bytes into a stream of control sequences,
character tokens, and category-code annotations, mirroring the behavior
described in \texttt{tex.web}. Section~\ref{sec:macro-expansion} discusses
how these tokens are subsequently expanded.

\subsubsection{Category Codes}

Every input character is assigned one of sixteen category codes before any
macro expansion takes place. This table is user-configurable via
\texttt{\textbackslash catcode}, and the default assignment is summarized in
Table~\ref{tab:catcodes}.

\begin{table}[ht]
  \centering
  \begin{tabular}{|c|l|}
    \hline
    Code & Meaning \\
    \hline
    0 & Escape character \\
    1 & Begin group \\
    2 & End group \\
    6 & Parameter \\
    \hline
  \end{tabular}
  \caption{Selected default category codes}
  \label{tab:catcodes}
\end{table}

\subsection{Macro Expansion}
\label{sec:macro-expansion}

Macro expansion proceeds left to right, repeatedly replacing expandable
tokens until only primitive and unexpandable tokens remain, as summarized
by equation~\eqref{eq:expansion}:

\begin{equation}
  \label{eq:expansion}
  T_{n+1} = \operatorname{expand}(T_n), \qquad T_\infty = \lim_{n \to \infty} T_n.
\end{equation}

We revisit this fixed-point view again in Section~\ref{sec:layout} when
discussing how expansion interacts with the layout stage.

\section*{A Note on Compatibility}
\addcontentsline{toc}{section}{A Note on Compatibility}

This unnumbered section is nonetheless listed in the table of contents via
an explicit \texttt{\textbackslash addcontentsline} call, which is a common
pattern in reference documentation for marking important asides without
disturbing the numbered section sequence.

\section{Layout}
\label{sec:layout}

\subsection{Box and Glue Model}

Layout consumes the expanded token stream and produces a tree of boxes
connected by glue, following the classic TeX box-and-glue model. The
badness of a line is computed from the total stretch or shrink required to
reach the target width; see Table~\ref{tab:stages} for where this fits in
the overall pipeline.

\subsection{Page Breaking}

Page breaking operates on the vertical list produced by paragraph
breaking, choosing break points that minimize total badness across the
document, consistent with the discussion in Section~\ref{sec:tokenization}.

\section{Shaping}
\label{sec:shaping}

\subsection{Font Selection}

Font selection resolves each glyph run to a concrete font file and glyph
index, using the same NFSS conventions as \LaTeX{}, cross-referenced back
to the category-code table in Table~\ref{tab:catcodes}.

\subsection{Kerning and Ligatures}

Kerning and ligature substitution are applied per shaping run before glyphs
are positioned, a process that further refines the raw box widths computed
in Section~\ref{sec:layout}.

\section{Emission}
\label{sec:emission}

The final stage serializes the positioned glyph runs into PDF content
streams. Further implementation notes are available at
\url{https://example.org/flashtex/emission-notes}, and the full API
reference is linked from \href{https://example.org/flashtex/api}{the API
documentation}.

\section{Appendix: Glossary}
\label{sec:appendix}

\begin{itemize}
  \item \textbf{Badness}: a numeric measure of how far a line or page
        deviates from its ideal width or height.
  \item \textbf{Glue}: a flexible space with a natural size, stretch, and
        shrink, as introduced in Section~\ref{sec:layout}.
  \item \textbf{NFSS}: the New Font Selection Scheme used by \LaTeX{},
        referenced in Section~\ref{sec:shaping}.
\end{itemize}

This concludes the architecture reference; refer back to
Figure~\ref{fig:pipeline} and Table~\ref{tab:stages} for the high-level
summary.

\end{document}
